% problem-type: truyen-song-1d-dalembert
% order: 10
% Nguon: De thi giua ky MAT2306 (PTDHR), HK2 2025-2026, K68, De so 1, Cau 2. (tu giai)

\section{Truyền sóng 1D --- d'Alembert}

% ---------------------------------------------------------------
\subsection*{Đề bài ví dụ}
% ---------------------------------------------------------------
Xét phương trình truyền sóng $u_{tt}=u_{xx}$ ($x\in\mathbb{R}_+$, $t>0$), với điều kiện biên
$u_x(0,t)=0$, và điều kiện ban đầu
\[
  u(x,0)=\begin{cases}
    x^2 & 0\le x\le 1,\\
    2-x & 1\le x\le 2,\\
    0 & \text{còn lại},
  \end{cases}
  \qquad
  u_t(x,0)=\chi_{(0,1)}(x).
\]
\textbf{(a)} Thác triển các hàm điều kiện ban đầu thành hàm trên toàn trục. Xác định sóng tiến, sóng lùi.

\textbf{(b)} Vẽ nghiệm $u$ tại các thời điểm $t=\tfrac12,\,1,\,\tfrac32$. Hàm $u(x,\tfrac12)$ không liên tục (gãy) tại đâu?

\emph{\small(Nguồn: Đề thi giữa kỳ MAT2306 --- PTĐHR, HK2 2025--2026, K68, Đề số 1, Câu 2. Lời giải dưới đây tự soạn.)}

% ---------------------------------------------------------------
\subsection*{Nhận ra dạng này khi nào?}
% ---------------------------------------------------------------
\begin{itemize}
  \item PDE \textbf{1 chiều không gian}: $u_{tt}=c^2 u_{xx}$, dữ liệu ban đầu $(g,h)$.
  \item Hoặc bài nhiều chiều nhưng dữ liệu chỉ phụ thuộc một biến $\Rightarrow$ hạ về 1D.
  \item Có \textbf{biên} tại $x=0$ (nửa trục $\mathbb{R}_+$) $\Rightarrow$ phải thác triển dữ liệu ra toàn trục trước.
\end{itemize}
\begin{meobox}
Nghiệm 1D luôn tách thành \emph{sóng tiến} (chạy phải) $+$ \emph{sóng lùi} (chạy trái).
Biên \textbf{Neumann} $u_x(0,t)=0\Rightarrow$ thác triển \textbf{chẵn}; biên \textbf{Dirichlet} $u(0,t)=0\Rightarrow$ thác triển \textbf{lẻ}.
\emph{(Cách thác triển tổng quát + công thức tường minh: Phụ lục C --- Thác triển chẵn/lẻ.)}
\end{meobox}

% ---------------------------------------------------------------
\subsection*{Công thức \& phương pháp}
% ---------------------------------------------------------------
% method: cong-thuc-dalembert-1d
\begin{dieukienbox}
Dùng khi: sóng \textbf{1D} $u_{tt}=u_{xx}$ trên $\mathbb{R}$ (hoặc sau khi đã hạ bài toán về 1D).
\end{dieukienbox}

\begin{nhobox}
\textbf{Nghiệm tổng quát = sóng tiến + sóng lùi:}
\[
  u(x,t)=F(x-t)+G(x+t),
\]
$F(x-t)$ là \emph{sóng tiến} (giữ nguyên hình, chạy sang phải vận tốc $1$), $G(x+t)$ là \emph{sóng lùi} (chạy sang trái).

\textbf{Công thức sóng tiến / sóng lùi} (từ dữ liệu $u(x,0)=g$, $u_t(x,0)=h$):
\[
  F(\xi)=\tfrac12 g(\xi)-\tfrac12\!\int_0^{\xi}\! h(s)\,ds,\qquad
  G(\xi)=\tfrac12 g(\xi)+\tfrac12\!\int_0^{\xi}\! h(s)\,ds.
\]

\textbf{Công thức d'Alembert} (thay $F,G$ vào, hằng số khử nhau):
\[
  u(x,t)=\tfrac12\big(g(x-t)+g(x+t)\big)+\tfrac12\int_{x-t}^{x+t} h(\xi)\,d\xi .
\]
\end{nhobox}

\begin{meobox}
Điểm $x$ tại thời điểm $t$ chỉ chịu ảnh hưởng dữ liệu trong đoạn $[x-t,\,x+t]$ (miền phụ thuộc 1D).
\end{meobox}


% ---------------------------------------------------------------
\subsection*{Đề bài \& bài giải}
% ---------------------------------------------------------------
\paragraph{Nhắc lại đề.} Phương trình $u_{tt}=u_{xx}$ trên nửa trục $x>0$ (vận tốc $c=1$), biên $u_x(0,t)=0$ (Neumann --- tường ``tự do''),
với dữ liệu ban đầu
\[
  g(x):=u(x,0)=\begin{cases}x^2 & 0\le x\le1,\\ 2-x & 1\le x\le2,\\ 0 & x\ge2,\end{cases}
  \qquad
  h(x):=u_t(x,0)=\chi_{(0,1)}(x).
\]

\paragraph{Bước 1 --- Thác triển chẵn dữ liệu ra toàn trục (câu a).}\leavevmode

\emph{Bước 1.1 --- Vì sao thác triển chẵn?} Bài đang trên \emph{nửa đường} $x>0$, công thức d'Alembert lại đòi hỏi dữ liệu trên \emph{toàn trục} $\mathbb{R}$. Mẹo chuẩn: kéo dài $g,h$ sang $x<0$ sao cho biên tự thoả mãn. Vì hàm \textbf{chẵn} có $\partial_x|_{x=0}=0$ (đồ thị đối xứng qua trục tung nên tiếp tuyến ngang), nên nếu nghiệm là hàm chẵn theo $x$ thì $u_x(0,t)=0$ \emph{tự động}. Do tính tuyến tính, dữ liệu chẵn $\Rightarrow$ nghiệm chẵn. Quy tắc tổng quát:
\begin{itemize}
\item Biên Dirichlet ($u=0$) $\Rightarrow$ thác triển \textbf{lẻ} (hàm lẻ triệt tiêu tại $0$).
\item Biên Neumann ($u_x=0$) $\Rightarrow$ thác triển \textbf{chẵn} (hàm chẵn có đạo hàm $=0$ tại $0$).
\end{itemize}

\emph{Bước 1.2 --- Công thức tường minh.} Thác triển chẵn $g_e(x)=g(|x|)$, $h_e(x)=h(|x|)$:
\[
  g_e(x)=\begin{cases}
    x^2 & |x|\le 1,\\
    2-|x| & 1\le |x|\le 2,\\
    0 & |x|\ge 2,
  \end{cases}
  \qquad
  h_e(x)=\chi_{(-1,1)}(x).
\]
(Trên $x<0$: $g_e(-x)=g(x)$; $h$ trở thành đặc trưng của khoảng đối xứng $(-1,1)$.)

\paragraph{Bước 2 --- Áp công thức d'Alembert toàn trục.}\leavevmode

\emph{Bước 2.1 --- Ráp công thức.} Trên $\mathbb{R}$ với dữ liệu $(g_e,h_e)$ và $c=1$, công thức d'Alembert cho
\[
  u(x,t)=\tfrac12\big(g_e(x-t)+g_e(x+t)\big)+\tfrac12\!\int_{x-t}^{x+t}\! h_e(s)\,ds.
\]
(Hệ số $\tfrac{1}{2c}=\tfrac12$ ở số hạng vận tốc do $c=1$.)

\emph{Bước 2.2 --- Tách sóng tiến / sóng lùi.} Đây là dạng $u=F(x-t)+G(x+t)$ với
\[
  F(\xi)=\tfrac12 g_e(\xi)-\tfrac12\!\int_0^{\xi}\! h_e(s)\,ds,
  \qquad
  G(\xi)=\tfrac12 g_e(\xi)+\tfrac12\!\int_0^{\xi}\! h_e(s)\,ds.
\]
\begin{itemize}
  \item $F(x-t)$ là \textbf{sóng tiến} (giữ nguyên hình dạng, chạy sang phải vận tốc $1$): trên đường đặc trưng $x-t=\text{const}$ giá trị $F$ không đổi.
  \item $G(x+t)$ là \textbf{sóng lùi} (chạy sang trái): trên đường đặc trưng $x+t=\text{const}$ giá trị $G$ không đổi.
\end{itemize}
Hai họ đường thẳng $\{x-t=\text{const}\}$ và $\{x+t=\text{const}\}$ trong mặt phẳng $(x,t)$ chính là \textbf{hai họ đặc trưng}. \textbf{Miền phụ thuộc} của điểm $(x,t)$ là đoạn $[x-t,x+t]$ trên trục ban đầu.

\emph{Bước 2.3 --- Ý nghĩa phản xạ ở biên.} Với $x>0$ và $t>x$, đặc trưng trái rơi vào miền âm: $x-t<0$. Khi đó tính chẵn cho
\(g_e(x-t)=g_e(t-x)\), tức sóng lẽ ra đi sang trái khỏi tường $x=0$ thì \emph{phản xạ ngược lại} và giữ nguyên dấu. Đây là cách biên Neumann được mã hoá ngầm qua thác triển chẵn.

\paragraph{Bước 3 --- Tính phần vận tốc.}
Vì $h_e=\chi_{(-1,1)}$, tích phân chính là độ đo Lebesgue của giao:
\[
  \tfrac12\!\int_{x-t}^{x+t}\! h_e(s)\,ds
  =\tfrac12\,\big|\,[x-t,\,x+t]\cap(-1,1)\,\big|.
\]
Hàm này \textbf{liên tục theo $x$} (tích phân hàm bị chặn là liên tục Lipschitz).

\paragraph{Bước 4 --- Câu (b): liệt kê điểm gãy của $u(\cdot,\tfrac12)$.}\leavevmode

\emph{Bước 4.1 --- Tìm các kỳ dị của dữ liệu thác triển.} Kiểm tra liên tục của $g_e$: tại $x=1$ giá trị $1=1$, tại $x=2$ giá trị $0=0$, tại $x=0$ thì $g_e(0^\pm)=0=0$ --- $g_e$ \emph{liên tục} khắp nơi nhưng \emph{đạo hàm nhảy} (góc gãy) tại
\[
  x=0\ (\text{độ dốc } -2\to +2),\quad
  x=\pm1\ (\text{từ } \pm2\ \text{sang} \mp1),\quad
  x=\pm2\ (\text{từ } \mp1\ \text{sang } 0).
\]
Còn $h_e=\chi_{(-1,1)}$ có \textbf{bước nhảy} tại $x=\pm1$ (và \emph{không} có bước nhảy tại $x=0$ vì $h_e(0^-)=h_e(0^+)=1$ --- đây là ưu điểm của thác triển chẵn).

\emph{Bước 4.2 --- Lan truyền kỳ dị dọc theo đặc trưng.} Trong công thức d'Alembert: mỗi kỳ dị của $g_e$ tại $x=x_0$ sinh ra góc gãy của $u(\cdot,t)$ tại $x_0\pm t$ (qua hai số hạng $g_e(x\mp t)$); mỗi bước nhảy của $h_e$ tại $x=x_1$ cũng sinh ra góc gãy của $u$ tại $x_1\pm t$ (qua tích phân; bước nhảy của tích phân biến thành góc gãy của $u$). Bản thân $u$ luôn \textbf{liên tục} (không có bước nhảy).

\emph{Bước 4.3 --- Áp dụng tại $t=\tfrac12$, chỉ giữ $x>0$.} Dịch các điểm $\{0,\pm1,\pm2\}$ (kỳ dị của $g_e$) và $\{\pm1\}$ (bước nhảy của $h_e$) đi $\pm\tfrac12$:
\[
  \{0,\pm1,\pm2\}\pm\tfrac12=\Big\{\pm\tfrac12,\pm\tfrac32,\pm\tfrac52,\pm\tfrac32,\pm\tfrac12\Big\},\qquad
  \{\pm1\}\pm\tfrac12=\Big\{\pm\tfrac12,\pm\tfrac32\Big\}.
\]
Lấy hợp rồi giữ $x>0$: ta có các ứng viên $x=\tfrac12,\tfrac32,\tfrac52$.

\begin{nhobox}
$u(x,\tfrac12)$ \textbf{liên tục} trên $[0,\infty)$ nhưng \textbf{không khả vi} (đạo hàm gãy) tại
\[
  x=\tfrac12,\quad x=\tfrac32,\quad x=\tfrac52.
\]
Ý ``không liên tục'' trong đề chỉ \textbf{đạo hàm} $u_x$ bị gián đoạn, còn bản thân $u$ không có bước nhảy.
\end{nhobox}
